\section{Completion of the positive comparison}\label{sec:positive-completion}

We complete the proof of Proposition~\ref{prop:positive}. All curves in
this section are non-CM. We first compute the central value of
\eqref{eq:L1} when the analytic rank is already one. Together with the
rank-zero calculation, this transfers a unit from the positive twist to
the prescribed curve.

For an initially unknown center of Selmer corank one, we then choose a
tame direction in which the Selmer Bockstein is nonzero. Along that
direction the unit determinant forces the singular localization of the
Kato class to have a simple zero. The torsion test in
Proposition~\ref{tame:spectral} converts its nonzero derivative into
analytic rank one. This argument uses a compact Selmer class before
knowing that it is generated by a rational point.

Isogeny invariance from Section~\ref{sec:classical} allows us to work
with \(E'\), transporting the positive anchor hypothesis to
\(E'^{(h)}\) when \(h\) is the given discriminant. Use the square orders,
true \(2\)-towers, symbols with \(A=1\), and fixed \(c,d\) of the
positive determinant construction. Include each moving prime once in
the orders, and include the fixed twist in the fixed support when it
is used. Thus Proposition~\ref{prop:siegel-norms} removes any moving
prime at its trivial-variable specialization, with the other
parameters still allowed. All auxiliary sequences retain the splitting
and trace bounds of \eqref{eq:L1}. The option to prescribe one of them
will supply the Bockstein direction.

\subsection{Central specialization in analytic rank one}

\begin{lemma}[The rank-one central value]\label{lem:positive-completion:rank-one}
Suppose a branch \(A_0\) of the positive determinant construction has analytic
rank one. Then \(U(0)\ne0\) and
\[
                         v_2 U(0)=X(A_0).
\]
Thus \eqref{eq:L5} holds in both analytic ranks zero and one.
\end{lemma}
\begin{proof}
\emph{The compact Selmer line.}
Let \(A_0\) be \(E'\) or the indicated positive twist. Work with the notation \(O^j,R^j_\mathbb R,J_2,\tau_g,\tau_q,b\) of Lemma~\ref{lem:positive-determinant:rank-zero}, at the actual stages (so \(S'=S\cup R_i\), \(\omega\) minimal on \(A_0\)). Write \(P\) for a rational point generating the free Mordell--Weil lattice, and \(s_1=v_2(\#\Sha(A_0/\mathbb Q))\), finite by the low analytic rank theorem.

The Poitou--Tate map from \(J_2\) to the dual of the discrete Selmer
group is injective. Indeed its composite with the dual of
\(A_0(\mathbb Q)\otimes\mathbb Q_2/\mathbb Z_2\) is evaluation by
compact point localization. Log duality and \(\log_\omega P\ne0\)
make this composite injective, with cokernel of length
\(v_2(\log_\omega P)-b\). The Kummer exact sequence adds the finite
quotient of length \(s_1\), so the cokernel into the dual Selmer has
length \(s_1+v_2(\log_\omega P)-b\).

The preceding term in Poitou--Tate is the singular image of ordinary
\(O^1\), which must therefore be zero. At the other places compact
\(H^1\) is already Kummer. Global Kummer and finiteness of \(\Sha\)
now identify \(O^1\) with the compact Mordell--Weil group, and the
remaining terms of Poitou--Tate give
\[
 \begin{split}
     O^1&= A_0(\mathbb Q)^\wedge_2,\\
     \operatorname{length} O^2&=s_1+v_2(\log_\omega P)-b+
                \operatorname{length}R^2_\mathbb R+
                         \sum_{q\in S'_f}\tau_q-\tau_g .
 \end{split}
\]

\emph{Removing the moving Euler factors.}
Write \(z^{\circ}_i\) for the ordinary image after specializing all \(B\)-parameters centrally on the given branch. The moving omitted factors there are \(e_r=1-a_{A_0}(r)/r+1/r\); the fixed twist in the comparison is split at these primes. Removing them all gives the fixed-support class \(z\) of the tame calculation, with
\(z^{\circ}_i=(\prod_{r\in R_i}e_r)z\) rationally after inflation, by the norm and tower relations. In particular by \eqref{eq:T1}
\[
 v_2\!\left(\frac{\log_\omega z^{\circ}_i}{(\log_\omega P)^2}\right)
  =v_2\!\left(\Delta(0)\,\prod_{q\in S'_f}L_q(A_0,1)^{-1}\cdot
              \frac{L'(A_0,1)}{\Omega_0 H(P)}\right).               \tag{P1}\label{eq:P1}
\]
Here \(\Delta(0)\) is the evaluated smoothing including any fixed quadratic scalar. To use \eqref{eq:T1} at the stripped base one makes \emph{its own} choice of horizontal sequence if necessary, killing the fixed support and fixed data required there and adapting to the point \(P\) on this actual curve; it need not be any current variable of \(B\).

\emph{The determinant valuation.}
The first equality in \eqref{eq:positive-determinant:raw-center} still applies, since \(O^1\) has a line, with nonzero \(z^\circ_i\) rationally by \eqref{eq:P1}, and \(O^2\) is torsion. Now \(\operatorname{length}(O^1/\mathbb Z_2 z^\circ_i)=v_2(\log_\omega z^\circ_i/\log_\omega P)+\tau_g\). Thus the raw valuation is
\[
 v_2\!\left(\frac{\log_\omega z^\circ_i}{(\log_\omega P)^2}\right)
            +b+2\tau_g-s_1-\sum_{q\in S'_f}\tau_q
                             -\operatorname{length}R^1_\mathbb R
       = X(A_0)+v_2(\Delta(0)),
\]
by the Haar and real-component computations. Here \({\rm Reg}_{A_0}=H(P)\) on the full free lattice. Just as at rank zero these are bounded-torsion-length central specializations: the moving \(\tau_r\) are bounded, since \(r\to1\) and the Tate trace limits are not 2 there. The displayed cohomology formulas and positive triangle bound positive torsion and keep the rational degree-one dimension two with other degrees zero. Thus pivots for those dimensions work at the center of the limit diagram also; the class determinant specializes there and has the same finite valuation. The fixed corrections have central ratio 1. This proves \eqref{eq:L5} in analytic rank one.
\end{proof}

\begin{corollary}[Positivity and transfer of a central unit]\label{cor:positive-completion:unit}
The following assertions hold for a non-CM curve \(A\).
\begin{enumerate}[label=(\roman*)]
\item If \(\mathrm{an}(A)\le1\), then \(X(A)\ge0\).
\item Suppose the positive twist hypothesis of Proposition~\ref{prop:positive}
holds, without assuming the analytic rank of \(A\). Then the positive
determinant of \(E'\) is a unit for every completed auxiliary-prime
construction permitted by Proposition~\ref{prop:positive-determinant:integral},
including those with one prescribed admissible sequence. If \(s_2(E')=0\),
then \(\mathrm{an}(E')=0\) and \(X(E')=0\).
\end{enumerate}
\end{corollary}
\begin{proof}
Integrality in Proposition~\ref{prop:positive-determinant:integral}, the
central equality \eqref{eq:L5} in both ranks, and isogeny invariance prove
the lower bound. Under the positive twist hypothesis, the central value of
\(U\) on that branch has valuation zero. Since \(B\) is local, \(U\) is
a unit on that branch. By the simultaneous construction and reduction-unit transfer for \eqref{eq:L1}, \(U\) at \(E'\) is then a unit as well. If the given twist is trivial we simply have the unit directly. This conclusion works for any of the completed choices in that construction, including with the prescribed-sequence option. Thus for \(s_2(E')=0\), Lemma~\ref{lem:positive-determinant:rank-zero} gives analytic
nonvanishing and \(X(E')=0\).
\end{proof}

\subsection{A direction for the unknown corank-one center}\label{subsec:positive-completion:direction}

Assume now \(s_2(E')=1\) and the positive anchor hypothesis. Write
\(E=E'\) for the untwisted branch, and use \(f,N,T\) as in
Section~\ref{sec:tame}. Two-primary parity gives odd functional sign;
the rational Mordell--Weil rank is still unknown.

Choose \(K,\mathcal A,\ell\) and the other fixed data of the tame
spectral-height comparison, including the nonvanishing companion.
Include the anchor's twisting field among the splitting constraints.
We will prescribe the first auxiliary sequence \(r_{1,i}=r_i\), with
residue character \(\lambda_i\), to satisfy Definition~\ref{tame:data}
and a nonzero Bockstein condition. In particular its character kills
all primes of \(S_f,c,d\). We establish this compatibility before
completing the auxiliary tuple.

\begin{lemma}[A primitive for the self-cup]\label{lem:positive-completion:self-cup}
Let \(x\) be a cocycle representing a nonzero class in the rational compact Selmer line at fixed
support, with coefficients \(T\otimes\mathbb Q_2\), and let
\(\beta=\operatorname{pol}(x)\) on the Tate dual. There is a continuous scalar
\(1\)-cochain \(k\) on \(G_{\mathbb Q,S}\) such that
\[
                              dk=-\beta\cup x.
\]
After multiplying \(x,\beta\) by the same nonzero dyadic integer and
\(k\) by its square, all three cochains can be taken integral.
\end{lemma}
\begin{proof}
The inverse limit of the finite Selmer groups fits into the Kummer
sequence with the rational-point completion and the Tate module of
\(\Sha\). Its rational dimension is therefore the corank \(s_2=1\)
of the usual divisible-coefficient Selmer group. These local
conditions may be imposed on \(G_{\mathbb Q,S}\), since at good odd
places outside \(S\) they are unramified.

The class of \(\beta\cup x\) vanishes locally by Kummer isotropy.
Moreover,
\(H^2(G_{\mathbb Q,S},\mathbb Q_2(1))\) injects into the sum of local
cohomology groups by Poitou--Tate: the opposite everywhere-locally-zero
kernel for \(\mathbb Q_2\) is unramified everywhere and vanishes by
class field theory. Thus the global cup is a coboundary, giving \(k\).
Compactness bounds the denominators of these continuous cochains.
Rescale \(x,\beta\) together and \(k\) quadratically to make them
integral. No rational-point representative of \(x\) is required.
\end{proof}

Fix a nonzero compact Selmer class and cochains \(x,\beta,k\) as in
Lemma~\ref{lem:positive-completion:self-cup}.

As in the choice of direction for \eqref{eq:T2}, let \(\mathcal P\) be the
finite list of rational primes at which the residue characters are required
to vanish, and set
\[
 L=F\bigl(\mu_{2^\infty},\,q^{1/2^n}:q\in\mathcal P,\ n\ge1\bigr).
\]
Here \(F/\mathbb Q\) is finite normal and enforces the fixed splitting and
congruence conditions, including the Hilbert class field, the fixed twist
data and \(E[2]\).

\begin{lemma}[A compatible nonzero direction]\label{lem:positive-completion:direction}
The prescribed auxiliary sequence can be chosen to satisfy the
spectral-height hypotheses of Section~\ref{sec:tame} and the fixed splitting
conditions for both positive-determinant branches. Its Frobenius lifts can
be chosen to tend on the fixed data to an element \(g\in G_L\) whose
Tate action \(g_T\) has determinant one, trace \(a_*\ne\pm2\), and two
non-root-of-unity eigenvalues, and which satisfies
\[
                   k(g)-\beta(g)g_T(g_T-1)^{-1}x(g)\ne0.             \tag{P2}\label{eq:P2}
\]
\end{lemma}
\begin{proof}
Apply Lemma~\ref{tame:cochain-direction} to the nonzero compact
Selmer class and its self-cup primitive from
Lemma~\ref{lem:positive-completion:self-cup}, with the displayed field
\(L\). Its hypotheses require only a nonzero cohomology class and
these fixed-support cochains. The prescribed list includes every
spectral-height condition and the fixed data of both branches.

The lemma gives \(g\) satisfying \eqref{eq:P2} and horizontal
Frobenius/residue-character data approaching it. Splitting on \(E[2]\)
and in the twisting field, together with the determinant-one trace
condition, also gives the prescribed-prime hypotheses of
Proposition~\ref{prop:positive-determinant:integral}. Write \(v=u_1\),
so the Galois scalar is \((1+v)^{\lambda_i}\), up to the consistent
inverse orientation.
\end{proof}

Complete \(R_i\) after this choice by Lemma~\ref{lem:positive-determinant:residual} (notably its mutual-bit dummy choice when the scalar-constituent test has two primes). It does not matter if \(\lambda_i\) fails to kill the supplementary \(r_{j,i}, j\ne1\). Their trace-limit and splitting conditions still hold. Corollary~\ref{cor:positive-completion:unit} therefore gives a unit \(U\)
on the \(E\)-branch with this prescribed variable included.

\subsection{The Selmer Bockstein is nonzero}\label{subsec:positive-completion:bockstein}

Pass from \(B\) to the germ \(\mathcal R=\mathbb Q_2[[v]]\) by
setting every other parameter to zero. Let \(C_{\rm glob}\) be the
ordinary limiting global complex on this germ. At \(2\) its coefficient
is constant because \(\lambda_i(2)=0\). We may therefore impose the
scalar extension of the rational Kummer line in degree one. Denote
the resulting Selmer complex by \(C_f^{\rm Sel}\), and the local
quotient by \(J\). There are no local invariant terms, so \(J\) is a
free singular line in degree one, measured by dual exponential, and
\[
                 C_f^{\rm Sel}\longrightarrow C_{\rm glob}
                                      \longrightarrow J
\]
is the local-condition triangle. Polarization identifies the use of
\(T\) with that of \(T^\vee(1)\), with the indicated scalar twist.

No other local conditions are needed. At the closed point, every odd
finite local complex in use is rationally acyclic. For a moving place,
the unramified local model has both \(g_j-1\) and
\(g_j/\lim_i r_{j,i}-1\) invertible. Unramified inflation thus
identifies closed-point ordinary cohomology with the fixed-support
problem. The model and inflation comparisons of
Section~\ref{sec:foundations}, after inverting \(2\), also give
perfectness over \(\mathcal R\).

At the closed point, exact Selmer orthogonality and Poitou--Tate give
one line in each of degrees one and two: the degree-one dimension is
\(s_2=1\), and global and local invariant vanishings exclude the
remaining degrees. A minimal model over \(\mathcal R\) consequently
has the form
\[
       [\,\mathcal R\xrightarrow{\delta(v)}\mathcal R\,]
       \quad\text{in degrees }1,2,
       \qquad \delta(0)=0.
\]
The first-order differential \(\delta'(0)\) is the Selmer Bockstein.
The next lemma proves that it is nonzero.

\begin{lemma}[Nonzero Selmer Bockstein]\label{lem:positive-completion:bockstein}
For the direction of Lemma~\ref{lem:positive-completion:direction},
\(\operatorname{ord}_v\delta(v)=1\). In particular,
\(C_f^{\rm Sel}\) is generically acyclic over \(\mathcal R\).
\end{lemma}
\begin{proof}
Suppose that \(\delta(v)\) has order at least two, including the
possibility that it vanishes identically. Then the central class
\(\beta\) lifts to the Selmer problem over \(\mathcal R/v^2\).
We first realize this hypothetical lift on finite-stage cochains,
then test it by cup reciprocity.

\emph{Realizing the first-order lift.}
After a fixed nonzero dyadic multiplier \(m\) and passage to a slower
cofinal precision, the lift gives cocycles on
\(G_{\mathbb Q,S\cup R_i}\) with the deformed coefficients,
\[
 b_i(v)=m\beta+v z_{1,i}\pmod{(2^{n_i},v^2)},
 \qquad n_i\le m_i,\qquad n_i\longrightarrow\infty,
\]
with the following two properties:
\begin{equation}\label{eq:positive-completion:stage-lift}
 d z_{1,i}=-m\lambda_i\cup\beta\pmod{2^{n_i}},
 \qquad
 \operatorname{inv}_2(z_{1,i}\cup x)\longrightarrow0.
\end{equation}
Use \(-\lambda_i\) instead throughout for the inverse scalar action.
These derivatives are consequences of the hypothetical lift; they
are not Kato classes.

Here are the finite-model details. After the one-variable
specialization, transfer restriction at \(2\) to constant local
models; the coefficients there are already split at each stage.
At dual-number precision retain the cup with the fixed \(x\), the
local invariant map on root coefficients, and the global model image
of the inflated central \(\beta\). A rational lift is a model
cocycle modulo \(v^2\) whose constant term represents \(\beta\).
Adjust it by a deformed coboundary to match that model image, then
clear one fixed common denominator.

The cycle and constant-term equations pass to stage representatives
modulo \((2^{n_i},v^2)\) at a slower cofinal precision, by the marked
finite-model construction. Include these representatives back into
actual cochains. The central contraction homotopy expresses the
difference of their constant terms from \(m\beta\) as a coboundary.
Lift the adjusting degree-zero cochain and subtract its deformed
coboundary. The constant term is now exactly \(m\beta\), so the
coefficient of \(v\) in the cocycle equation gives the first identity
in \eqref{eq:positive-completion:stage-lift}.

On the constant local models at \(2\), the derivative of the
hypothetical Selmer lift is rationally finite. Its cup invariant
against \(x\) is therefore zero by Kummer isotropy. The marked cup
and invariant maps make the stage invariants converge to this value.
The last coboundary adjustment does not change them because
\(\lambda_i=0\) locally at \(2\). This proves the second assertion
in \eqref{eq:positive-completion:stage-lift} using actual finite-stage
cochains.

\emph{The closed scalar cocycle.}
Suppress the index on \(z_{1,i}\). The cochain
\[
                   z_1\cup x+m\lambda_i\cup k
                   \tag{P3}\label{eq:P3}
\]
is closed: its two differentials cancel by
\(dz_1=-m\lambda_i\cup\beta\) and \(dk=-\beta\cup x\).
At \(2\), the second term is zero and the first has invariant tending
to zero by \eqref{eq:positive-completion:stage-lift}.
At the other fixed places, \(\lambda_i\) vanishes and the rational
\(x\) is a coboundary, so one common multiplier kills their local
contributions. Outside the allowed support the cup is unramified.

\emph{The moving local invariants.}
At every moving prime there is a bounded-denominator local cochain
\(w_{j,i}\) with \(dw_{j,i}=x\): since \(x\) is unramified, take
\[
 w_{j,i}=(\rho_T(\operatorname{Fr}_{j,i})-1)^{-1}
                                      x(\operatorname{Fr}_{j,i}).
\]
The trace bounds give the common denominator. After clearing it,
\eqref{eq:P3} is locally cohomologous to
\(m\lambda_i\cup(k-\beta\cup w_{j,i})\). The second factor is an
unramified root-coefficient cocycle. At a supplementary prime the
character \(\lambda_i\) is also unramified, hence
\[
 \operatorname{inv}_{r_{j,i}}\bigl(m\lambda_i\cup
                   (k-\beta\cup w_{j,i})\bigr)=0
                    \qquad(j\ne1).
\]

At the distinguished prime \(r_i\), local cup reciprocity and the
residue-exponent choice instead give
\[
 \operatorname{inv}_{r_i}\bigl(m\lambda_i\cup
                   (k-\beta\cup w_{1,i})\bigr)
       \longrightarrow \pm m\eta\ne0,
\]
with the same fixed denominator clearing understood. To see the
normalization, at the working precision the local field contains
the relevant roots of unity. An unramified root cocycle is the Kummer
class of a unit, and its Frobenius value records raising the residue
unit to \((r_i-1)/2^{n_i}\). This is exactly the additive residue
exponent defining \(\lambda_i\), in the compatible root basis.
Frobenius convergence gives the scalar \(\eta\) of \eqref{eq:P2}.

All other local invariants tend to zero. The nonzero limit at \(r_i\)
contradicts the global sum-of-invariants law for the closed cocycle
\eqref{eq:P3}. Thus \(\delta(v)\) has order one.
\end{proof}

\subsection{Detection of the analytic zero}

\begin{proof}[Completion of the proof of Proposition~\ref{prop:positive}]
The inequality and the corank-zero implication were proved in
Corollary~\ref{cor:positive-completion:unit}. Suppose \(s_2(E)=1\), and use
the direction and the unit constructed above.
By Lemma~\ref{lem:positive-completion:bockstein}, the local-condition
triangle gives a single generic degree-one line for the ordinary
global problem on this germ. Restoring the ordinary real triangle
gives just degree one of dimension two for the positive complex.
These are the same dimensions as on the full ring \(B\), so the
generic determinant identity defining \(U\) specializes to
\(\operatorname{Frac}\mathcal R\), changing pivots if necessary.

The real basis generates the constant real contribution, and every
correction factor in \eqref{eq:L1} is a unit in the germ. Write
\(z^{\rm all}\) for the Kato class with all supplementary moving
primes still present. The local-condition triangle and the inverse
determinant convention now give
\begin{equation}\label{eq:positive-completion:simple-exponential}
 \begin{split}
 0=\operatorname{ord}_v U
   &=d_{\mathcal R}(C_f^{\rm Sel},1)
       +d_{\mathcal R}(J,\operatorname{loc}_{\rm sing}z^{\rm all})\\
   &=-1+\operatorname{ord}_v(\exp^*z^{\rm all}).
 \end{split}
\end{equation}
Thus the dual exponential has a simple zero, in any fixed nonzero
rational differential coordinate.

Remove the supplementary primes at each stage by the good Euler norm,
with the true cyclotomic variable trivial and the conductor-\(r_i\)
variable retained. The Euler polynomials restoring those primes may
have nontrivial unramified scalar values, but their limiting central
values are \(2-\operatorname{tr}g_j\ne0\). They are therefore units in
\(\mathcal R\). The rational norm identities pass to the local scalar
polynomials and their limits: the Euler factors are integral, and the
split local dual-exponential coordinates have uniformly bounded
denominators. Consequently the one-prime horizontal comparison also
has a simple zero, and its derivative \(Z'\) is nonzero, up to a
possible variable inversion.

If \(L'(f,1)=0\), the nonzero companion value and \eqref{eq:GZ}
make \(\phi_*P_X\) torsion. Proposition~\ref{tame:spectral} would then
give \(Z'=0\), a contradiction. Its torsion implication requires no
rational spanning point, and the prescribed direction satisfies all
its field, eigenvalue, and residue-character conditions. Since the
functional sign is odd, we conclude that \(L'(E,1)\ne0\) and
\(\mathrm{an}(E)=1\).

We may now apply Lemma~\ref{lem:positive-completion:rank-one}.
The unit \(U\) gives \(X(E)=0\); the classical low-rank theorem gives
the corresponding rational rank and finiteness of \(\Sha\). Isogeny
invariance transports these conclusions from \(E'\) to the original
curve.
\end{proof}  

This proves Proposition~\ref{prop:positive}. We next establish the split-pair
anchor of Proposition~\ref{prop:anchor} and the CM comparison of
Proposition~\ref{prop:cm}.

